\chapter{Convergence and the Central Limit Theorem}

\begin{center}\itshape
Roll a hundred dice and add them up. The total is $350$ on average, but what is the chance that it exceeds\\
$400$? A hundred dice have $6^{100}$ outcomes, and yet the answer, about $0.15\%$, comes from one function,\\
the bell curve, which knows nothing about dice.
\end{center}

In 1733 Abraham de Moivre, a French refugee who earned his living in London coffee houses solving gamblers' problems, found that the number of heads
in many tosses of a coin follows a curve shaped like a bell. Laplace extended the discovery to other games; Lyapunov, in 1901, to almost every sum of
independent quantities. Francis Galton was so struck by it that he wrote of the bell curve: ``It reigns with serenity and in complete self-effacement,
amidst the wildest confusion. The huger the mob, and the greater the apparent anarchy, the more perfect is its sway.'' This chapter proves it.

The law of large numbers says that an average of $n$ independent copies of $X$ is close to $\mathbb E[X]$. The central
limit theorem says how close, and in which way: the error, magnified by $\sqrt n$, is approximately Gaussian, whatever
the law of $X$, as soon as its variance is finite. This universality is why the bell curve appears in measurement
errors, in heights, in polls and in the noise of electronic circuits.

The chapter first sorts out what ``converge'' means for random variables: almost surely, in probability, in mean,
or in law. It proves the strong law of large numbers under a fourth moment, and its consequence for empirical
distribution functions. Then it proves the central limit theorem by Lindeberg's replacement method, which only
uses Taylor's formula: the variables are swapped one at a time for Gaussian ones, and each swap costs little.
Characteristic functions, the Fourier transforms of laws, appear as a second tool, strong enough to show that a law
is determined by its Fourier transform and that averaging Cauchy variables is useless. The chapter ends with the
local version of de Moivre and Laplace, with Stirling's constant computed from the central limit theorem, and with
Gaussian vectors, the multivariate bell curves promised in Volume~0.

\section{Modes of convergence}

\begin{definition}
Let $X,X_1,X_2,\ldots$ be random variables. We say that $X_n\to X$
\begin{itemize}
\item \emph{almost surely} if $P(X_n\to X)=1$;
\item \emph{in probability} if $P(|X_n-X|>\varepsilon)\to0$ for every $\varepsilon>0$;
\item \emph{in $L^p$} ($p=1,2$) if $\mathbb E|X_n-X|^p\to0$;
\item \emph{in law} if $\mathbb E[f(X_n)]\to\mathbb E[f(X)]$ for every bounded continuous $f:\R\to\R$.
\end{itemize}
\end{definition}

The first three need all the variables on one probability space; convergence in law only concerns their laws, and
we also write $X_n\to\mu$ in law when $\mu$ is the law of $X$. The weak law of Chapter~4 is a convergence in probability,
and Corollary~\ref{cor:slln} an almost sure one.

\begin{think}
Let $A_1,A_2,\ldots$ be independent events with $P(A_n)=\frac1n$. Does $\mathbf 1_{A_n}$ converge to $0$? In which of the four senses?
\end{think}

\begin{proposition}\label{prop:modes}
\begin{enumerate}[label=(\alph*)]
\item Almost sure convergence and convergence in $L^p$ each imply convergence in probability.
\item Convergence in probability implies convergence in law; and if $X_n\to c$ in law for a constant $c$, then $X_n\to c$ in
probability.
\item If $X_n\to X$ in probability, a subsequence converges almost surely.
\item No other implication holds in general: there are sequences that converge in probability and in $L^p$ but not almost
surely; almost surely but not in $L^1$; and in law but not in probability.
\end{enumerate}
\end{proposition}
\begin{proof}
(a) If $X_n\to X$ a.s., then $\mathbf 1_{|X_n-X|>\varepsilon}\to0$ a.s., and dominated convergence (bounded by $1$) gives $P(|X_n-X|>\varepsilon)\to0$. In $L^p$,
Markov's inequality gives $P(|X_n-X|>\varepsilon)\le\varepsilon^{-p}\mathbb E|X_n-X|^p$.

(b) Let $f$ be bounded by $K$ and continuous, and $\varepsilon\in\,]0,1]$. Choose $M$ with $P(|X|>M)<\varepsilon$ (continuity of $P$ along $\{|X|>M\}\downarrow\varnothing$), and
by Heine (Lemma~\ref{lem:heine}) on $[-M-1,M+1]$, a $\delta\in\,]0,1]$ with $|f(x)-f(y)|<\varepsilon$ when $x,y$ lie there and $|x-y|<\delta$. If $|X|\le M$ and
$|X_n-X|<\delta$, then $|f(X_n)-f(X)|<\varepsilon$; otherwise it is at most $2K$. So $\mathbb E|f(X_n)-f(X)|\le\varepsilon+2K\varepsilon+2KP(|X_n-X|\ge\delta)$, and the last term
tends to $0$. For the converse with a constant, let $f(x)=\min(1,|x-c|/\varepsilon)$: then $P(|X_n-c|\ge\varepsilon)\le\mathbb E[f(X_n)]\to f(c)=0$.

(c) Choose $n_1<n_2<\cdots$ with $P(|X_{n_k}-X|>2^{-k})<2^{-k}$. By Borel--Cantelli (Theorem~\ref{thm:bc}(a)), almost surely $|X_{n_k}-X|\le2^{-k}$ for all
large $k$.

(d) With the $A_n$ of the think box, $X_n=\mathbf 1_{A_n}$ satisfies $P(|X_n|>\varepsilon)\le\mathbb E|X_n|^p=\frac1n\to0$, but $X_n=1$ infinitely often almost surely by
Borel--Cantelli (b). With $U$ uniform on $]0,1[$, $X_n=n\mathbf 1_{U<1/n}\to0$ everywhere, but $\mathbb E[X_n]=1$. With $X$ a $N(0,1)$ variable, $X_n=-X$ has
the law of $X$ for every $n$, so $X_n\to X$ in law, but $|X_n-X|=2|X|$ does not tend to $0$ in probability.
\end{proof}

Convergence in law can be read on distribution functions, and can be tested with smooth functions only. A
non-decreasing $F$ has at most countably many jumps (at most $k$ jumps of size $\ge\frac1k$), so its continuity points are dense.

\begin{proposition}\label{prop:law}
Let $F_n$ and $F$ be the distribution functions of $X_n$ and $X$.
\begin{enumerate}[label=(\alph*)]
\item $X_n\to X$ in law if and only if $F_n(t)\to F(t)$ at every point $t$ where $F$ is continuous.
\item If $\mathbb E[f(X_n)]\to\mathbb E[f(X)]$ for every $f$ of class $C^3$ with $f,f',f'',f'''$ bounded, then $X_n\to X$ in law.
\end{enumerate}
\end{proposition}
\begin{proof}
Let $t$ be a continuity point of $F$ and $\delta>0$. Suppose $g\le\mathbf 1_{]-\infty,t]}\le h$ for functions $g$, $h$ with values in $[0,1]$ such that $g=1$ on
$]-\infty,t-\delta]$ and $h=0$ on $[t+\delta,\infty[$. If $\mathbb E[g(X_n)]\to\mathbb E[g(X)]$ and likewise for $h$, then
\[
\begin{aligned}F(t-\delta)\le\mathbb E[g(X)]&\le\liminf F_n(t)\eqbr\le\limsup F_n(t)\\&\le\mathbb E[h(X)]\le F(t+\delta),\end{aligned}
\]
and letting $\delta\to0$ gives $F_n(t)\to F(t)$. For (a), take $g$ and $h$ piecewise linear. For (b), let $H(u)=1$ for $u\le0$, $H(u)=0$ for $u\ge1$ for $0\le u\le1$,
\[
H(u)=1-u^4(35-84u+70u^2-20u^3).
\]
Then $H(1)=0$ and $H'(u)=-140u^3(1-u)^3$ on $[0,1]$, so $H$ is non-increasing and its first three derivatives vanish at $0$ and $1$: $H$ is of
class $C^3$ on $\R$, with bounded derivatives. Take $g(x)=H(\frac{x-t+\delta}\delta)$ and $h(x)=H(\frac{x-t}\delta)$. This proves (b) once the converse half of (a) is
established.

Conversely, suppose $F_n\to F$ at continuity points, and let $f$ be continuous with $|f|\le K$, and $\varepsilon>0$. Choose continuity points $a<b$
of $F$ with $F(a)<\varepsilon$ and $1-F(b)<\varepsilon$, and by Heine continuity points $a=s_0<s_1<\dots<s_m=b$ such that $f$ varies by less than $\varepsilon$ on each
$[s_{j-1},s_j]$. Let $g=\sum_jf(s_j)\mathbf 1_{]s_{j-1},s_j]}$. Then $|f-g|\le\varepsilon$ on $]a,b]$ and $|f-g|\le K$ elsewhere, so for $Y=X_n$ or $X$,
$|\mathbb E f(Y)-\mathbb E g(Y)|\le\varepsilon+K\,P(Y\notin\,]a,b])$. Now $\mathbb E g(X_n)=\sum_jf(s_j)(F_n(s_j)-F_n(s_{j-1}))\to\mathbb E g(X)$ and
$P(X_n\notin\,]a,b])=F_n(a)+1-F_n(b)\to F(a)+1-F(b)<2\varepsilon$. Hence $\limsup|\mathbb E f(X_n)-\mathbb E f(X)|\le2\varepsilon+4K\varepsilon$.
\end{proof}

\section{The strong law of large numbers}

\begin{theorem}[strong law, fourth moment]\label{thm:slln4}
Let $X_1,X_2,\ldots$ be independent with the same mean $\mu$ and $\sup_i\mathbb E[(X_i-\mu)^4]=K<\infty$. Then $\frac{S_n}n\to\mu$ almost surely and in $L^2$.
\end{theorem}
\begin{proof}
Let $Y_i=X_i-\mu$, so $\mathbb E Y_i=0$. Expand $\mathbb E[(Y_1+\dots+Y_n)^4]$ as a sum of $\mathbb E[Y_iY_jY_kY_l]$. By independence, a term in which some index
appears exactly once vanishes, since it factors with a factor $\mathbb E[Y_i]=0$. What remains is the $n$ terms $\mathbb E[Y_i^4]$ and, for each of the
$\binom n2$ pairs $i<j$, the $6$ arrangements of $Y_i^2Y_j^2$, with $\mathbb E[Y_i^2Y_j^2]=\mathbb E[Y_i^2]\mathbb E[Y_j^2]\le\sqrt{\mathbb E[Y_i^4]\,\mathbb E[Y_j^4]}\le K$ by Cauchy--Schwarz.
So
\[
\mathbb E\bigl[(S_n-n\mu)^4\bigr]\le nK+3n(n-1)K\le3n^2K,
\]
and by Markov, $P(|\frac{S_n}n-\mu|\ge\varepsilon)\le\frac{3K}{n^2\varepsilon^4}$, which is summable. Conclude as in Corollary~\ref{cor:slln} with Borel--Cantelli. The $L^2$
convergence is the weak law's computation: $\mathbb E[(\frac{S_n}n-\mu)^2]\le\frac{\sqrt K}n$, since $\operatorname{Var}X_i\le\sqrt K$.
\end{proof}

The strong law holds for every i.i.d.\ sequence with $\mathbb E|X_1|<\infty$ (Kolmogorov 1933; Etemadi's 1981 proof uses only
pairwise independence), with a truncation argument that we leave for later in the series. Its most used consequence
is that simulation works: the empirical distribution of a sample converges to the true law.

\begin{theorem}[Glivenko--Cantelli]\label{thm:gc}
Let $X_1,X_2,\ldots$ be i.i.d.\ with a continuous distribution function $F$, and $F_n(t)=\frac1n\#\{i\le n:X_i\le t\}$ the empirical distribution
function. Then $\sup_{t\in\R}|F_n(t)-F(t)|\to0$ almost surely.
\end{theorem}
\begin{proof}
Fix $k\ge1$. By the intermediate value theorem there are $t_1\le\dots\le t_{k-1}$ with $F(t_j)=\frac jk$; put $t_0=-\infty$, $t_k=+\infty$ (with $F_n=F=0$ at
$-\infty$ and $=1$ at $+\infty$). For each $j$, $F_n(t_j)$ is the average of the independent bounded variables $\mathbf 1_{X_i\le t_j}$ of mean $F(t_j)$, so by
Corollary~\ref{cor:slln} the event $E_k=\{F_n(t_j)\to F(t_j)\text{ for all }j\}$ has probability $1$. For $t_{j-1}\le t<t_j$, monotonicity gives
\[
\begin{aligned}F_n(t)-F(t)&\le F_n(t_j)-F(t_{j-1})\eqbr=F_n(t_j)-F(t_j)+\tfrac1k,\\
F_n(t)-F(t)&\ge F_n(t_{j-1})-F(t_j)\eqbr=F_n(t_{j-1})-F(t_{j-1})-\tfrac1k.\end{aligned}
\]
So $\sup_t|F_n(t)-F(t)|\le\max_j|F_n(t_j)-F(t_j)|+\frac1k$, and on $E_k$ its $\limsup$ is at most $\frac1k$. On $\bigcap_kE_k$, an event of probability $1$, it tends to $0$.
\end{proof}

Kolmogorov showed that $\sqrt n\sup_t|F_n-F|$ has a
limit law that does not depend on $F$; this is the Kolmogorov--Smirnov test used to check random number generators,
as numerically of Chapter~5. The argument above uses continuity of $F$ to choose the quantiles $t_j$, and the case of a jump is not proved here.

\section{Characteristic functions}

For a random vector $X$ in $\R^d$, the \emph{characteristic function} is $\varphi_X(t)=\mathbb E\bigl[e^{i\langle t,X\rangle}\bigr]=\mathbb E\cos\langle t,X\rangle+i\,\mathbb E\sin\langle t,X\rangle$
for $t\in\R^d$, the Fourier transform of the law. It always exists, $|\varphi_X|\le1$ (by Jensen for the convex modulus, or by
Cauchy--Schwarz on the real and imaginary parts), $\varphi_X(0)=1$, $\varphi_X$ is continuous by dominated convergence, and
$\varphi_{aX+b}(t)=e^{itb}\varphi_X(at)$ for real $X$. If $X$ and $Y$ are independent, $\varphi_{X+Y}=\varphi_X\varphi_Y$: expand $e^{i\langle t,X\rangle}e^{i\langle t,Y\rangle}$ into real and
imaginary parts and use $\mathbb E[UV]=\mathbb E[U]\mathbb E[V]$ for bounded independent $U$, $V$ (Proposition~\ref{prop:indepcdf}).

\begin{lemma}[differentiation under the integral]\label{lem:diffint}
Let $g(t,x)$ be defined for $t$ in an open interval $J$ and $x\in E$, integrable in $x$ for each $t$, and differentiable in $t$ with
$|\partial_tg(t,x)|\le G(x)$ for an integrable $G$. Then $t\mapsto\int g(t,x)\dd\mu(x)$ is differentiable on $J$, with derivative $\int\partial_tg(t,x)\dd\mu(x)$.
\end{lemma}
\begin{proof}
For $h_k\to0$, the quotients $\frac{g(t+h_k,x)-g(t,x)}{h_k}$ equal $\partial_tg(t+\theta h_k,x)$ for some $\theta\in\,]0,1[$ (mean value theorem), so they are bounded by $G(x)$,
and they tend to $\partial_tg(t,x)$. Apply dominated convergence.
\end{proof}

\begin{proposition}\label{prop:cfexamples}
\begin{enumerate}[label=(\alph*)]
\item If $X$ is $N(m,\sigma^2)$, then $\varphi_X(t)=e^{imt-\sigma^2t^2/2}$.
\item If $X$ has the Laplace density $\frac12e^{-|x|}$, then $\varphi_X(t)=\frac1{1+t^2}$.
\end{enumerate}
\end{proposition}
\begin{proof}
(a) By the scaling rule it suffices to take $m=0$, $\sigma=1$. Since the density $\phi(x)=\frac1{\sqrt{2\pi}}e^{-x^2/2}$ is even, $\mathbb E\sin(tX)=0$ and $\varphi(t)=\int\cos(tx)\phi(x)\dd x$.
By the lemma with $G(x)=|x|\phi(x)$, $\varphi'(t)=-\int x\sin(tx)\phi(x)\dd x$. Since $-x\phi(x)=\phi'(x)$, an integration by parts (on $[-A,A]$, then $A\to\infty$) gives
$\varphi'(t)=\int\sin(tx)\phi'(x)\dd x=-t\int\cos(tx)\phi(x)\dd x=-t\varphi(t)$. So $(\varphi(t)e^{t^2/2})'=0$, and $\varphi(t)=\varphi(0)e^{-t^2/2}=e^{-t^2/2}$.
(b) By symmetry $\varphi(t)=\int_0^\infty e^{-x}\cos(tx)\dd x=:I$. Two integrations by parts give $I=1-t\int_0^\infty e^{-x}\sin(tx)\dd x=1-t^2I$.
\end{proof}

\begin{theorem}[uniqueness and inversion]\label{thm:inversion}
\begin{enumerate}[label=(\alph*)]
\item Two random vectors in $\R^d$ with the same characteristic function have the same law.
\item If $X$ is real and $\int_\R|\varphi_X(t)|\dd t<\infty$, then $X$ has the continuous density $f(x)=\frac1{2\pi}\int_\R e^{-itx}\varphi_X(t)\dd t$.
\end{enumerate}
\end{theorem}
\begin{proof}
Let $Z$ be a vector of $d$ independent $N(0,1)$ variables, independent of $X$ (by (A1) we may enlarge the probability space,
which does not change the law of $X$), and $\sigma>0$. The density of $\sigma Z$ is $\psi_\sigma(y)=\prod_j\frac1{\sigma\sqrt{2\pi}}e^{-y_j^2/(2\sigma^2)}$, and by (a) of
Proposition~\ref{prop:cfexamples} applied to a $N(0,\sigma^{-2})$ variable, coordinate by coordinate,
\[
\psi_\sigma(y)=\frac1{(2\pi)^d}\int_{\R^d}e^{i\langle t,y\rangle}e^{-\sigma^2|t|^2/2}\dd t.
\]
By Tonelli and the translation invariance of $\lambda_d$ (the one-dimensional case in each coordinate), $X+\sigma Z$ has the density
$p_\sigma(x)=\mathbb E[\psi_\sigma(x-X)]$. Insert the formula for $\psi_\sigma$ and exchange $\mathbb E$ and $\int\dd t$ (Fubini: the integrand has modulus
$e^{-\sigma^2|t|^2/2}$, integrable on $\Omega\times\R^d$):
\[
p_\sigma(x)=\frac1{(2\pi)^d}\int_{\R^d}e^{i\langle t,x\rangle}\varphi_X(-t)\,e^{-\sigma^2|t|^2/2}\dd t.
\]
(a) So the law of $X+\sigma Z$ depends only on $\varphi_X$. For bounded continuous $g$, $\mathbb E[g(X+\sigma Z)]\to\mathbb E[g(X)]$ as $\sigma\to0$ by dominated
convergence; so $\mathbb E[g(X)]$ is determined by $\varphi_X$ for every such $g$. For a box $B=\prod_j\,]-\infty,s_j]$, the continuous functions
$g_k(x)=\prod_j\min\bigl(1,\max(0,1-k(x_j-s_j))\bigr)$ decrease to $\mathbf 1_B$, so $P(X\in B)=\lim\mathbb E[g_k(X)]$ is determined too. Boxes form a
$\pi$-system generating $\mathcal B(\R^d)$: the law is determined (Theorem~\ref{thm:uniqueness}).

(b) Here $d=1$. Substituting $t\to-t$,
\[
p_\sigma(x)=\frac1{2\pi}\int e^{-itx}\varphi_X(t)e^{-\sigma^2t^2/2}\dd t,
\]
and
\[
\begin{aligned}|p_\sigma(x)-f(x)|&\eqbr\le\frac1{2\pi}\int|\varphi_X(t)|\bigl(1-e^{-\sigma^2t^2/2}\bigr)\dd t,\end{aligned}
\]
which tends to $0$ as $\sigma\to0$, uniformly in $x$, by dominated convergence. So $f$ is real, non-negative, and continuous (dominated convergence again). For a
continuous $g$ vanishing outside a segment, $\mathbb E[g(X)]=\lim_{\sigma\to0}\int gp_\sigma=\int gf$. Trapezoidal $g_k$ equal to $1$ on $[a,b]$ and $0$ outside
$[a-\frac1k,b+\frac1k]$ decrease to $\mathbf 1_{[a,b]}$, so $P(a\le X\le b)=\int_a^bf$ for all $a<b$, and $f$ is a density of $X$.
\end{proof}

\begin{corollary}[averaging Cauchy variables]\label{cor:cauchy}
A Cauchy variable has characteristic function $e^{-|t|}$. If $X_1,\dots,X_n$ are independent Cauchy variables, then $\frac{X_1+\dots+X_n}n$ is
again Cauchy.
\end{corollary}
\begin{proof}
Apply Theorem~\ref{thm:inversion}(b) to the Laplace law, whose characteristic function $\frac1{1+t^2}$ is integrable:
$\frac12e^{-|x|}=\frac1{2\pi}\int e^{-itx}\frac{\dd t}{1+t^2}$, that is, $\int e^{-itx}\frac{\dd t}{\pi(1+t^2)}=e^{-|x|}$. Read with $x=-s$, this says that the Cauchy law has
characteristic function $e^{-|s|}$. The average then has characteristic function $(e^{-|t|/n})^n=e^{-|t|}$, and (a) concludes.
\end{proof}

The average of a million Cauchy variables is as spread out as a single one: the law of large numbers fails, as it
must without a mean (Exercise~\ref{ex:cauchy}). Lévy's continuity theorem, that $X_n\to X$ in law iff $\varphi_{X_n}\to\varphi_X$ pointwise, gives
another proof of the central limit theorem; it is proved in Volume~10, together with the Berry--Esseen bound below.

\section{The central limit theorem}

\begin{think}
Throw $100$ fair dice. The total has mean $350$ and standard deviation $\sqrt{100\cdot\frac{35}{12}}\approx17.1$. Estimate the probability that it
exceeds $400$, using only the bell curve and the fact that $400$ is about $2.9$ standard deviations above the mean.
\end{think}

Let $\Phi(t)=\int_{-\infty}^t\frac1{\sqrt{2\pi}}e^{-x^2/2}\dd x$ be the $N(0,1)$ distribution function. The theorem is stated for a \emph{triangular array}: for
each $n$, a finite family of independent variables, possibly with different laws.

\begin{theorem}[Lindeberg, 1922]\label{thm:lindeberg}
For each $n$ let $X_{n,1},\dots,X_{n,k_n}$ be independent, with $\mathbb E[X_{n,k}]=0$, variances $\sigma_{n,k}^2$ with $\sum_k\sigma_{n,k}^2=1$, and suppose that for every
$\varepsilon>0$
\[
L_n(\varepsilon)=\sum_{k=1}^{k_n}\mathbb E\bigl[X_{n,k}^2\mathbf 1_{|X_{n,k}|>\varepsilon}\bigr]\xrightarrow[n\to\infty]{}0.
\]
Then $T_n=X_{n,1}+\dots+X_{n,k_n}$ converges in law to $N(0,1)$: $P(T_n\le t)\to\Phi(t)$ for every $t$.
\end{theorem}
\begin{proof}
Fix $f$ of class $C^3$ with $f''$ and $f'''$ bounded, and write $R(x,h)=f(x+h)-f(x)-f'(x)h-\frac12f''(x)h^2$. By Taylor's formula to order $3$,
$|R(x,h)|\le\frac16\|f'''\|\,|h|^3$, and to order $2$, $f(x+h)-f(x)-f'(x)h=\frac12f''(\xi)h^2$, so $|R(x,h)|\le\|f''\|h^2$. Hence, for any $\varepsilon>0$,
\[
|R(x,h)|\le\tfrac16\|f'''\|\,\varepsilon h^2+\|f''\|h^2\mathbf 1_{|h|>\varepsilon}.\tag{$*$}
\]
Drop the variables with $\sigma_{n,k}=0$, which vanish almost surely. Let $Z_{n,k}$ be $N(0,\sigma_{n,k}^2)$, independent of each other and of the $X_{n,k}$
(enlarging the probability space by (A1)). By Proposition~\ref{prop:gausssum}, $N=\sum_kZ_{n,k}$ is $N(0,1)$. Fix $n$, drop it from the notation,
and let $H_k=\sum_{j<k}Z_j+\sum_{j>k}X_j$ be the \emph{hybrid} in which the first $k-1$ variables have been replaced. Since $H_k+Z_k=H_{k+1}+X_{k+1}$,
the sum telescopes:
\[
f(T_n)-f(N)=\sum_{k=1}^{k_n}\bigl(f(H_k+X_k)-f(H_k+Z_k)\bigr).
\]
Now $H_k$ is independent of $X_k$ and of $Z_k$ (it is a function of the other variables), and $X_k$, $Z_k$ have the same first two moments. So
\[
\begin{aligned}\mathbb E[f(H_k+X_k)]&\eqbr=\mathbb E[f(H_k)]+\mathbb E[f'(H_k)]\,\mathbb E[X_k]\eqbr+\tfrac12\mathbb E[f''(H_k)]\,\mathbb E[X_k^2]\eqbr+\mathbb E[R(H_k,X_k)],\end{aligned}
\]
and likewise with $Z_k$; the first three terms are the same, and $|\mathbb E f(H_k+X_k)-\mathbb E f(H_k+Z_k)|\le\mathbb E|R(H_k,X_k)|+\mathbb E|R(H_k,Z_k)|$. By $(*)$,
$\sum_k\mathbb E|R(H_k,X_k)|\le\frac16\|f'''\|\varepsilon+\|f''\|L_n(\varepsilon)$. For the Gaussian terms, $\mathbb E|Z_k|^3=\sigma_k^3\,\mathbb E|N|^3\le2\sigma_k^3$ (since $\mathbb E|N|^3=2\sqrt{2/\pi}$), so
\[
\begin{aligned}\sum_k\mathbb E|R(H_k,Z_k)|&\le\tfrac13\|f'''\|\sum_k\sigma_k^3\eqbr\le\tfrac13\|f'''\|\max_k\sigma_k,\end{aligned}
\]
and $\sigma_k^2=\mathbb E[X_k^2]\le\varepsilon^2+\mathbb E[X_k^2\mathbf 1_{|X_k|>\varepsilon}]\le\varepsilon^2+L_n(\varepsilon)$, so $\max_k\sigma_k\le\sqrt{\varepsilon^2+L_n(\varepsilon)}$. Altogether
\[
\bigl|\mathbb E f(T_n)-\mathbb E f(N)\bigr|\le\tfrac16\|f'''\|\varepsilon+\|f''\|L_n(\varepsilon)+\tfrac13\|f'''\|\sqrt{\varepsilon^2+L_n(\varepsilon)}.
\]
As $n\to\infty$, $L_n(\varepsilon)\to0$, and therefore
\[
\limsup_n\bigl|\mathbb E f(T_n)-\mathbb E f(N)\bigr|\le\tfrac16\|f'''\|\varepsilon+\tfrac13\|f'''\|\varepsilon
\]
for every $\varepsilon>0$. The right side may be taken arbitrarily small, so $\mathbb E f(T_n)\to\mathbb E f(N)$. By Proposition~\ref{prop:law}(b) and the continuity of $\Phi$, $P(T_n\le t)\to\Phi(t)$ for all $t$.
\end{proof}

\begin{figure}[ht]
\centering
\fig{ch06-lindeberg}
\caption{Lindeberg's replacement. Each row is a hybrid sum: from top to bottom, the variables $X_k$ (orange) are replaced one at a
time by Gaussian variables $Z_k$ (blue) with the same mean and variance. The top row is $T_n$, the bottom row is exactly $N(0,1)$, and
each step changes $\mathbb E f$ only by a third-order Taylor remainder.}
\label{fig:lindeberg}
\end{figure}

\begin{corollary}[central limit theorem]\label{cor:clt}
Let $X_1,X_2,\ldots$ be i.i.d.\ with mean $\mu$ and variance $\sigma^2\in\,]0,\infty[$, and $S_n=X_1+\dots+X_n$. Then for every $t$,
\[
P\Bigl(\frac{S_n-n\mu}{\sigma\sqrt n}\le t\Bigr)\xrightarrow[n\to\infty]{}\Phi(t).
\]
\end{corollary}
\begin{proof}
Take $X_{n,k}=\frac{X_k-\mu}{\sigma\sqrt n}$ for $k\le n$. Then $\sum_k\mathbb E[X_{n,k}^2]=1$ and
\[
L_n(\varepsilon)=\mathbb E\Bigl[\frac{(X_1-\mu)^2}{\sigma^2}\mathbf 1_{|X_1-\mu|>\varepsilon\sigma\sqrt n}\Bigr],
\]
which tends to $0$ by
dominated convergence (dominated by $\frac{(X_1-\mu)^2}{\sigma^2}$, integrable).
\end{proof}

For the hundred dice, $P(S>400)=P(S\ge401)\approx1-\Phi(\frac{400.5-350}{17.08})=1-\Phi(2.96)\approx0.00155$, where $400.5$ rather than $400$ is the usual
correction for an integer-valued sum; the exact value, computed numerically, is $0.00151$. For Bernoulli
variables the corollary is the theorem of de Moivre (1733, for $p=\frac12$) and Laplace (1812): the binomial law is
approximately Gaussian.

\begin{figure}[ht]
\centering
\fig{ch06-sums}
\caption{The density of $\frac{S_n-n}{\sqrt n}$, where $S_n$ is a sum of $n$ independent exponential$(1)$ variables (a Gamma$(n)$ law), for $n=1$, $2$, $4$ and
$16$, against the $N(0,1)$ density (black). The skewness of the exponential law fades like $\frac1{\sqrt n}$.}
\label{fig:sums}
\end{figure}

\begin{remark}
How fast? The Berry--Esseen theorem says that
\[
\begin{aligned}\sup_t\Bigl|P\Bigl(\frac{S_n-n\mu}{\sigma\sqrt n}\le t\Bigr)-\Phi(t)\Bigr|&\eqbr\le\frac{C\,\mathbb E|X_1-\mu|^3}{\sigma^3\sqrt n},\end{aligned}
\]
with an absolute constant $C$
known to be less than $0.48$ (Shevtsova, 2011). The rate $\frac1{\sqrt n}$ is attained for coin flips, because of the jumps of the binomial
distribution function. It is proved in Volume~10.
\end{remark}

\begin{beauty}
Lindeberg's proof never looks at the law of $X$ beyond its first two moments: replacing one summand by a
Gaussian of the same mean and variance changes $\mathbb E f$ by a third-order term, and $n$ such changes of size $n^{-3/2}$
add up to almost nothing. The same \emph{swapping} idea proves universality theorems far beyond this chapter, for the
eigenvalues of random matrices (Tao and Vu) and for the outcomes of Boolean functions (the invariance principle of
Mossel, O'Donnell and Oleszkiewicz).
\end{beauty}

\section{The local theorem and Stirling's constant}

The central limit theorem is about $P(S_n\le t)$. For coin flips one can say more: each individual probability $P(S_n=k)$ is
close to the bell curve, and the proof is Stirling's formula together with the exponent $I_p$ of Chapter~4.

\begin{theorem}[local de Moivre--Laplace]\label{thm:local}
Let $S_n$ be binomial$(n,p)$ with $0<p<1$, $q=1-p$, $\phi(x)=\frac1{\sqrt{2\pi}}e^{-x^2/2}$ and $x_k=\frac{k-np}{\sqrt{npq}}$. For every $A>0$,
\[
\sup_{k:\,|x_k|\le A}\ \Bigl|\frac{\sqrt{npq}\;P(S_n=k)}{\phi(x_k)}-1\Bigr|\xrightarrow[n\to\infty]{}0.
\]
\end{theorem}
\begin{proof}
Write $m!=\sqrt{2\pi m}\,(\frac me)^me^{\varepsilon_m}$, with $\varepsilon_m\to0$ (Volume~0, Theorems~5.6 and~6.4). Let $a=\frac kn$; then $|a-p|\le\frac{A\sqrt{pq}}{\sqrt n}$, so $k$ and
$n-k$ tend to infinity uniformly. As in the proof of Theorem~\ref{thm:cramer},
\[
\begin{aligned}P(S_n=k)&=\binom nkp^kq^{n-k}\eqbr=e^{\varepsilon_n-\varepsilon_k-\varepsilon_{n-k}}\frac1{\sqrt{2\pi na(1-a)}}\,e^{-nI_p(a)}.\end{aligned}
\]
Uniformly in the allowed $k$: $a(1-a)\to pq$, the factor $e^{\varepsilon_n-\varepsilon_k-\varepsilon_{n-k}}$ tends to $1$, and it remains to see that $nI_p(a)-\frac{x_k^2}2\to0$.
Now $I_p(p)=I_p'(p)=0$, $I_p''(a)=\frac1{a(1-a)}$ so $I_p''(p)=\frac1{pq}$, and $I_p'''(a)=\frac1{(1-a)^2}-\frac1{a^2}$ is bounded by some $C$ for $a$ in $[\frac p2,\frac{1+p}2]$. By
Taylor's formula, $|I_p(a)-\frac{(a-p)^2}{2pq}|\le\frac C6|a-p|^3$, and since $\frac{x_k^2}2=\frac{n(a-p)^2}{2pq}$,
$|nI_p(a)-\frac{x_k^2}2|\le\frac{Cn}6|a-p|^3\le\frac{C(A\sqrt{pq})^3}{6\sqrt n}\to0$.
\end{proof}

\begin{figure}[ht]
\centering
\fig{ch06-local}
\caption{The binomial$(40,0.3)$ probabilities (bars) and $\frac1{\sqrt{npq}}\phi(x_k)$ (curve). Even at $n=40$ the agreement is close; it is worst in the
tails, where the theorem says nothing, and the relative error there is governed by $I_p$ rather than by its quadratic
approximation.}
\label{fig:local}
\end{figure}

Volume~0 proved Stirling's formula in two steps: $n!\sim C\sqrt n(\frac ne)^n$ for some constant $C$, then $C=\sqrt{2\pi}$ by Wallis' integrals.
The central limit theorem finds the constant independently, since the only $\sqrt{2\pi}$ it uses is the Gaussian
integral.

\begin{theorem}[Stirling's constant from the CLT]\label{thm:stirlingclt}
$\dfrac{n!\,e^n}{n^{n+1/2}}\to\sqrt{2\pi}$.
\end{theorem}
\begin{proof}
Let $S_n$ be Poisson$(n)$, the law of a sum of $n$ independent Poisson$(1)$ variables (Chapter~3: PGFs multiply), with mean and
variance $1$ each. By Corollary~\ref{cor:clt}, $Y_n=\frac{S_n-n}{\sqrt n}\to N$ in law, where $N$ is $N(0,1)$. First, an exact computation: since
\[
\sum_{k=0}^n(n-k)\frac{n^k}{k!}=\sum_{k=0}^n\frac{n^{k+1}}{k!}-\sum_{k=1}^n\frac{n^k}{(k-1)!}=\frac{n^{n+1}}{n!}
\]
(the second sum is the first without its last term), $\mathbb E[(n-S_n)^+]=e^{-n}\frac{n^{n+1}}{n!}$, that is, $\mathbb E[g(Y_n)]=\frac{n^{n+1/2}e^{-n}}{n!}$ with $g(y)=\max(-y,0)$.
Second, $\mathbb E[g(Y_n)]\to\mathbb E[g(N)]=\int_0^\infty x\phi(x)\dd x=\frac1{\sqrt{2\pi}}$. Indeed, $g_M=\min(g,M)$ is bounded and continuous, so $\mathbb E g_M(Y_n)\to\mathbb E g_M(N)$
(Proposition~\ref{prop:law}(a)); and $0\le g-g_M\le|y|\mathbf 1_{|y|>M}\le\frac{y^2}M$, with $\mathbb E[Y_n^2]=\mathbb E[N^2]=1$, so both errors are at most $\frac1M$.
\end{proof}

\section{Gaussian vectors}

\begin{think}
Let $X$ be $N(0,1)$ and $\varepsilon=\pm1$ a fair sign independent of $X$, and $Y=\varepsilon X$. Is $Y$ Gaussian? Are $X$ and $Y$ uncorrelated? Are they
independent?
\end{think}

\begin{definition}
A random vector $X=(X_1,\dots,X_d)$ is \emph{Gaussian} if $\langle u,X\rangle=\sum_ju_jX_j$ is Gaussian for every $u\in\R^d$, where constants count as
Gaussian $N(m,0)$.
\end{definition}

Each $X_j$ is then Gaussian, hence square-integrable, so $X$ has a mean $m=\mathbb E[X]$ and a \emph{covariance matrix} $\Sigma$ with
$\Sigma_{jk}=\operatorname{Cov}(X_j,X_k)$. By bilinearity $\operatorname{Var}\langle u,X\rangle=u^{\mathsf T}\Sigma u$, so $\Sigma$ is symmetric and positive semi-definite, and $\langle u,X\rangle$ is $N(\langle u,m\rangle,u^{\mathsf T}\Sigma u)$.

\begin{theorem}\label{thm:gaussvec}
\begin{enumerate}[label=(\alph*)]
\item A Gaussian vector with mean $m$ and covariance $\Sigma$ has characteristic function $e^{i\langle u,m\rangle-u^{\mathsf T}\Sigma u/2}$, so its law depends only on
$(m,\Sigma)$.
\item For every $m\in\R^d$ and every symmetric positive semi-definite $\Sigma$, there is a Gaussian vector with mean $m$ and covariance $\Sigma$.
\item The components of a Gaussian vector are independent if and only if they are uncorrelated ($\Sigma$ diagonal).
\item If $X$ is Gaussian, so is $AX+b$, with mean $Am+b$ and covariance $A\Sigma A^{\mathsf T}$. In particular, if $Z$ has independent $N(0,1)$
components and $Q$ is orthogonal, then $QZ$ has the law of $Z$.
\end{enumerate}
\end{theorem}
\begin{proof}
(a) $\varphi_X(u)$ is the characteristic function of the real variable $\langle u,X\rangle$ at $1$, given by Proposition~\ref{prop:cfexamples}(a) (directly if the
variance is $0$); then Theorem~\ref{thm:inversion}(a).
(b) By the spectral theorem (Volume~0, Theorem~7.7), $\Sigma=PDP^{\mathsf T}$ with $P$ orthogonal and $D$ diagonal with entries $\lambda_j\ge0$. Let $A=PD^{1/2}$ and
$X=m+AZ$ with $Z$ as in (d). Then $\langle u,X\rangle=\langle u,m\rangle+\sum_j(A^{\mathsf T}u)_jZ_j$ is a sum of independent Gaussian variables, hence Gaussian with
variance $|A^{\mathsf T}u|^2=u^{\mathsf T}AA^{\mathsf T}u=u^{\mathsf T}\Sigma u$ (Proposition~\ref{prop:gausssum}).
(c) Independent square-integrable variables are uncorrelated. Conversely, if $\Sigma$ is diagonal, the vector $(m_j+\sqrt{\Sigma_{jj}}Z_j)_j$ has independent
components and is Gaussian with the same mean and covariance, so by (a) it has the law of $X$; independence is a property
of the law.
(d) $\langle u,AX+b\rangle=\langle A^{\mathsf T}u,X\rangle+\langle u,b\rangle$ is Gaussian, and the mean and covariance follow by linearity. For $QZ$, the covariance is $QQ^{\mathsf T}=I$.
\end{proof}

In the think box, $Y$ is $N(0,1)$ (symmetry) and $\operatorname{Cov}(X,Y)=\mathbb E[\varepsilon]\,\mathbb E[X^2]=0$, but $|Y|=|X|$, so they are not independent: $(X,Y)$ is not a
Gaussian vector ($X+Y$ is $0$ with probability $\frac12$). Part (c) needs the \emph{joint} Gaussianity. When $\Sigma$ is invertible, a Gaussian
vector has the density $\frac1{(2\pi)^{d/2}\sqrt{\det\Sigma}}e^{-(x-m)^{\mathsf T}\Sigma^{-1}(x-m)/2}$, whose level sets are the ellipsoids of Volume~0's spectral theorem; this
needs the change of variables formula in $\R^d$ (Volume~6). The multivariate central limit theorem follows from
Corollary~\ref{cor:clt} applied to every $\langle u,X\rangle$, together with Lévy's theorem in $\R^d$ (Volume~10): sums of i.i.d.\ random vectors with a
covariance matrix are approximately Gaussian vectors.

\section{Exercises}

\begin{exercise}[$\star\star$]\label{ex:slutsky}
\emph{Slutsky's lemma.} Suppose $X_n\to X$ in law and $Y_n\to c$ in probability, for a constant $c$. Show that $X_n+Y_n\to X+c$ and
$X_nY_n\to cX$ in law.
\end{exercise}

\begin{exercise}[$\star$]\label{ex:confidence}
\emph{Monte Carlo confidence intervals.} Let $X_1,X_2,\ldots$ be i.i.d.\ with mean $\mu$ and variance $\sigma^2>0$, $\bar X_n$ the sample mean and
$\hat\sigma_n^2=\frac1n\sum_i(X_i-\bar X_n)^2$. Show that $P\bigl(|\bar X_n-\mu|\le1.96\,\hat\sigma_n/\sqrt n\bigr)\to2\Phi(1.96)-1\approx0.95$.
\end{exercise}

\begin{exercise}[$\star\star$]\label{ex:delta}
\emph{The delta method.} With $X_i$ as in Exercise~\ref{ex:confidence} and $g$ of class $C^1$ near $\mu$ with $g'(\mu)\ne0$, show that
$\sqrt n\bigl(g(\bar X_n)-g(\mu)\bigr)\to N(0,g'(\mu)^2\sigma^2)$ in law. Apply it to the estimate $4\bar X_n$ of $\pi$ by points in a quarter disc, and to $1/\bar X_n$.
\end{exercise}

\begin{exercise}[$\star$]\label{ex:cauchyclt}
Let $X_i$ be i.i.d.\ Cauchy. Show that $\frac{S_n}{\sqrt n}$ does not converge in law. Which hypothesis of the central limit theorem fails?
\end{exercise}

\begin{exercise}[$\star\star$]\label{ex:poissonhalf}
Show that $e^{-n}\sum_{k=0}^n\frac{n^k}{k!}\to\frac12$. (Ramanujan found the next term: the sum is $\frac12+\frac2{3\sqrt{2\pi n}}+O(\frac1n)$.)
\end{exercise}

\begin{exercise}[$\star\star$]\label{ex:returns}
Let $S_{2n}$ be the position after $2n$ steps of a simple random walk ($\pm1$ with probability $\frac12$). Show that $P(S_{2n}=0)\sim\frac1{\sqrt{\pi n}}$, and that
the expected number of returns to $0$ before time $2n$ is of order $\sqrt n$. (Chapter~7 builds on this.)
\end{exercise}

\begin{exercise}[$\star\star\star$]\label{ex:cyclesclt}
Let $C_n$ be the number of cycles of a uniformly random permutation of $\{1,\dots,n\}$. Show that $\frac{C_n-\ln n}{\sqrt{\ln n}}\to N(0,1)$ in law
(Goncharov, 1944). (Use Theorem~\ref{thm:feller} and Theorem~\ref{thm:lindeberg}.)
\end{exercise}

\section{Hints}

\begin{hint}{ex:slutsky}
By Proposition~\ref{prop:law}(b) it suffices to test smooth $f$ with bounded derivatives; compare $f(X_n+Y_n)$ with $f(X_n+c)$ using the mean value
theorem on $\{|Y_n-c|\le\delta\}$. For the product, first bound $P(|X_n|>M)$ uniformly in $n$.
\end{hint}
\begin{hint}{ex:confidence}
Show that $\hat\sigma_n\to\sigma$ almost surely (strong law, or the weak law for $X_i^2$ if you only know $\mathbb E X^2<\infty$: in probability suffices), and use
Slutsky.
\end{hint}
\begin{hint}{ex:delta}
$g(\bar X_n)-g(\mu)=g'(\xi_n)(\bar X_n-\mu)$ with $\xi_n$ between $\mu$ and $\bar X_n$.
\end{hint}
\begin{hint}{ex:cauchyclt}
$\frac{S_n}{\sqrt n}$ has the law of $\sqrt n\,X_1$; compute $P(|\sqrt nX_1|\le A)$.
\end{hint}
\begin{hint}{ex:poissonhalf}
The sum is $P(S_n\le n)$ for $S_n$ Poisson$(n)$.
\end{hint}
\begin{hint}{ex:returns}
$P(S_{2n}=0)=\binom{2n}n4^{-n}$; use Theorem~\ref{thm:local} with $k=n$, or Stirling directly.
\end{hint}
\begin{hint}{ex:cyclesclt}
$X_{n,j}=\frac{B_j-1/j}{s_n}$ with $s_n^2=\operatorname{Var}C_n\sim\ln n$; the variables are bounded by $\frac1{s_n}$. Then replace $s_n$ by $\sqrt{\ln n}$ and $H_n$ by $\ln n$ with
Slutsky.
\end{hint}

\section{Solutions}

\begin{solution}{ex:slutsky}
\idea A perturbation that is small in probability moves $f(X_n+\cdot)$ by little, except on an event of small probability.
\writeup Let $f$ be $C^3$ with bounded derivatives, $|f|\le K$, $|f'|\le K'$. On $\{|Y_n-c|\le\delta\}$, $|f(X_n+Y_n)-f(X_n+c)|\le K'\delta$; elsewhere it is at most $2K$.
So $|\mathbb E f(X_n+Y_n)-\mathbb E f(X_n+c)|\le K'\delta+2KP(|Y_n-c|>\delta)$, and $\mathbb E f(X_n+c)\to\mathbb E f(X+c)$ since $x\mapsto f(x+c)$ is bounded continuous. Let $n\to\infty$,
then $\delta\to0$; Proposition~\ref{prop:law}(b) concludes. For the product, choose $M$ with $P(|X|\ge M)<\varepsilon$ and $\pm M$ continuity points of $F_X$; by
Proposition~\ref{prop:law}(a), $P(|X_n|>M)<2\varepsilon$ for large $n$. On $\{|X_n|\le M,\ |Y_n-c|\le\delta\}$, $|f(X_nY_n)-f(cX_n)|\le K'M\delta$. So
$|\mathbb E f(X_nY_n)-\mathbb E f(cX_n)|\le K'M\delta+2K(2\varepsilon+P(|Y_n-c|>\delta))$, and $\mathbb E f(cX_n)\to\mathbb E f(cX)$. Let $n\to\infty$, then $\delta\to0$, then $\varepsilon\to0$.
\end{solution}

\begin{solution}{ex:confidence}
\idea The CLT holds with the true $\sigma$; replacing it by a consistent estimate changes nothing in the limit.
\writeup $\hat\sigma_n^2=\frac1n\sum X_i^2-\bar X_n^2$. Chebyshev's proof of the weak law (Theorem~\ref{thm:weaklaw}) works word for word with the integral of
Chapter~5, so $\bar X_n\to\mu$ in probability. The $Y_i=X_i^2$ are i.i.d.\ and integrable with mean $m=\sigma^2+\mu^2$, but may have no variance; truncate. For
$M>0$ let $Y_i'=Y_i\mathbf 1_{Y_i\le M}$ and $Y_i''=Y_i-Y_i'$, and $\eta_M=\mathbb E[Y_1'']$, which tends to $0$ as $M\to\infty$ by dominated convergence. Then
$\frac1n\sum_iY_i-m=\bigl(\frac1n\sum_iY_i'-\mathbb E Y_1'\bigr)+\frac1n\sum_iY_i''-\eta_M$. Chebyshev bounds the first bracket ($\operatorname{Var}Y_1'\le M^2$) and Markov the second, so
\[
P\Bigl(\Bigl|\frac1n\sum_iY_i-m\Bigr|>3\varepsilon\Bigr)\le\frac{M^2}{n\varepsilon^2}+\frac{\eta_M}\varepsilon
\]
as soon as $\eta_M\le\varepsilon$. Choose $M$ with $\eta_M\le\varepsilon^2$, then $n$ large: the average of the $Y_i$ tends to $m$ in probability. Hence $\hat\sigma_n^2\to\sigma^2$ and $\hat\sigma_n\to\sigma$ in probability, and $\frac\sigma{\hat\sigma_n}\to1$ in probability (on
$\{\hat\sigma_n>\sigma/2\}$, $|\frac\sigma{\hat\sigma_n}-1|\le\frac2\sigma|\hat\sigma_n-\sigma|$). By the CLT and Slutsky, $\frac{\sqrt n(\bar X_n-\mu)}{\hat\sigma_n}\to N(0,1)$, and since $\Phi$ is continuous,
$P(|\cdot|\le1.96)\to\Phi(1.96)-\Phi(-1.96)=2\Phi(1.96)-1$.
\end{solution}

\begin{solution}{ex:delta}
\idea Near $\mu$, $g$ is linear with slope $g'(\mu)$, and a linear map of an asymptotically Gaussian error is Gaussian.
\writeup $\sqrt n(g(\bar X_n)-g(\mu))=g'(\xi_n)\cdot\sqrt n(\bar X_n-\mu)$ (mean value theorem, on the event that $\bar X_n$ lies in the neighbourhood where $g$ is
$C^1$, whose probability tends to $1$). Since $\bar X_n\to\mu$ in probability, so does $\xi_n$, and $g'(\xi_n)\to g'(\mu)$ in probability by continuity of $g'$. The
second factor tends in law to $N(0,\sigma^2)$, so by Slutsky the product tends to $g'(\mu)N(0,\sigma^2)=N(0,g'(\mu)^2\sigma^2)$. For the quarter disc, $X_i$ is
Bernoulli$(\frac\pi4)$ and $g(x)=4x$: the error of $4\bar X_n$ is about $\frac4{\sqrt n}\sqrt{\frac\pi4(1-\frac\pi4)}\approx\frac{1.64}{\sqrt n}$. For $g(x)=\frac1x$ (with $\mu\ne0$), the variance is
$\frac{\sigma^2}{\mu^4}$.
\end{solution}

\begin{solution}{ex:cauchyclt}
\idea A Cauchy sum grows like $n$, not like $\sqrt n$.
\writeup By Corollary~\ref{cor:cauchy}, $\frac{S_n}n$ is Cauchy, so $\frac{S_n}{\sqrt n}=\sqrt n\frac{S_n}n$ has the law of $\sqrt nX_1$, and
$P(|\sqrt nX_1|\le A)=\frac2\pi\arctan\frac A{\sqrt n}\to0$ for every $A$. If $\frac{S_n}{\sqrt n}$ converged in law to some $Y$, then taking $\pm A$ continuity points with
$P(|Y|<A)>\frac12$ would give $P(|\frac{S_n}{\sqrt n}|\le A)\to P(-A\le Y\le A)>\frac12$ by Proposition~\ref{prop:law}(a), a contradiction. The variance is infinite.
\end{solution}

\begin{solution}{ex:poissonhalf}
\idea Read the sum as a probability for a sum of $n$ independent Poisson$(1)$ variables, and apply the CLT at $t=0$.
\writeup With $S_n$ Poisson$(n)$ as in Theorem~\ref{thm:stirlingclt}, the sum is $P(S_n\le n)=P(\frac{S_n-n}{\sqrt n}\le0)\to\Phi(0)=\frac12$.
\end{solution}

\begin{solution}{ex:returns}
\idea Returning to $0$ at time $2n$ is the central term of a binomial law.
\writeup $S_{2n}=0$ iff exactly $n$ of the $2n$ steps are $+1$: $P(S_{2n}=0)=\binom{2n}n4^{-n}$. Theorem~\ref{thm:local} with $2n$ trials, $p=\frac12$, $k=n$ and $x_k=0$ gives
$P(S_{2n}=0)\sim\frac{\phi(0)}{\sqrt{2n/4}}=\frac1{\sqrt{2\pi}}\sqrt{\frac2n}=\frac1{\sqrt{\pi n}}$. The expected number of returns at times $2,4,\dots,2n$ is
$\sum_{j\le n}P(S_{2j}=0)$; since the terms are equivalent to $\frac1{\sqrt{\pi j}}$ and $\sum_{j\le n}\frac1{\sqrt j}\sim2\sqrt n$ (compare with $\int_0^n\frac{\dd x}{\sqrt x}$), the
sum is $\sim2\sqrt{n/\pi}$.
\end{solution}

\begin{solution}{ex:cyclesclt}
\idea The number of cycles is a sum of independent Bernoulli$(\frac1j)$ variables; they are bounded and their total
variance $\ln n$ tends to infinity, which is all Lindeberg needs.
\writeup By Theorem~\ref{thm:feller}, $C_n$ has the law of $B_1+\dots+B_n$ with $B_j$ independent Bernoulli$(\frac1j)$, so we may assume equality. Let
$s_n^2=\operatorname{Var}C_n=H_n-\sum_jj^{-2}$, so $s_n^2\sim\ln n$ (as $H_n=\ln n+O(1)$ and $\sum j^{-2}\le\frac{\pi^2}6$). Put $X_{n,j}=\frac{B_j-1/j}{s_n}$: independent, centred, with
$\sum_j\mathbb E X_{n,j}^2=1$, and $|X_{n,j}|\le\frac1{s_n}$. For fixed $\varepsilon>0$, once $\frac1{s_n}<\varepsilon$ every indicator in $L_n(\varepsilon)$ vanishes, so $L_n(\varepsilon)=0$. By Lindeberg,
$\frac{C_n-H_n}{s_n}\to N(0,1)$. Finally
\[
\begin{aligned}\frac{C_n-\ln n}{\sqrt{\ln n}}&=\frac{s_n}{\sqrt{\ln n}}\cdot\frac{C_n-H_n}{s_n}\eqbr+\frac{H_n-\ln n}{\sqrt{\ln n}},\end{aligned}
\]
where $\frac{s_n}{\sqrt{\ln n}}\to1$ and the last term tends to $0$; Slutsky concludes. A random permutation of a million elements has on average
$H_{10^6}\approx14.4$ cycles, give or take $s_{10^6}\approx3.6$.
\end{solution}

\section*{Sources and further reading}

\begin{itemize}
\item J.~W.~Lindeberg, \emph{Eine neue Herleitung des Exponentialgesetzes in der Wahrscheinlichkeitsrechnung},
Mathematische Zeitschrift 15 (1922), 211--225. The replacement method, in its original form.
\item T.~Tao, \emph{Topics in Random Matrix Theory}, AMS, 2012, Section~2.2: the central limit theorem by several
methods, including Lindeberg's swapping and its modern descendants.
\item W.~Feller, \emph{An Introduction to Probability Theory and Its Applications}, Vol.~1, 3rd ed., Wiley, 1968,
Chapter~VII: de Moivre--Laplace, local and global.
\item R.~Durrett, \emph{Probability: Theory and Examples}, 5th ed., Cambridge University Press, 2019, Chapters~2--3:
the strong law, characteristic functions, Lévy's continuity theorem and Berry--Esseen.
\item I.~Shevtsova, \emph{On the absolute constants in the Berry--Esseen type inequalities for identically distributed
summands}, arXiv:1111.6554 (2011).
\item S.~M.~Stigler, \emph{The History of Statistics}, Harvard University Press, 1986. How de Moivre, Laplace and Gauss
found the bell curve.
\item V.~Goncharov, \emph{Du domaine d'analyse combinatoire}, Izvestiya Akad.\ Nauk SSSR 8 (1944), 3--48. Cycles of random
permutations and the CLT.
\item F.~Galton, \emph{Natural Inheritance}, Macmillan, 1889, p.~66, for the quotation; S.~M.~Stigler, \emph{The History of Statistics}, Harvard University Press, 1986, for de Moivre (1733), Laplace and the bell curve.
\end{itemize}
